\documentclass{article}

\textwidth=6.5in
\textheight=8.5in
\voffset=-54pt
\oddsidemargin=0in
\evensidemargin=0in
\setlength{\parskip}{8pt}
\setlength{\parindent}{0pt}

\usepackage{amsmath, amssymb}
\usepackage{ifthen}

\usepackage[inline]{enumitem}
\setlist{topsep=0pt, parsep=8pt}

\usepackage[rgb,table]{xcolor}\convertcolorsUfalse
\usepackage{graphicx}

% change hyperref options
\PassOptionsToPackage{hyphens}{url}
\usepackage[bookmarks,colorlinks,linkcolor=black,citecolor=blue,pdfstartview=FitH,urlcolor=blue]{hyperref}
\hypersetup{pdfpagemode=UseNone}
\usepackage{bookmark}

\usepackage{graphicx}
\usepackage{multicol}
\usepackage{framed}
\usepackage{array}

\usepackage{icomma}

\usepackage{pifont}
\newenvironment{noanswer}{}{}
\newlist{altenumerate}{enumerate}{3}

\usepackage{soul}

\usepackage{pgfplots}
\pgfplotsset{compat=1.18}  
\pgfplotscreateplotcyclelist{mycolor}{black, blue, red, green!70!black, magenta, purple, cyan, orange }  
\pgfplotsset{
  disabledatascaling,
  cycle list=tikzcolor, 
  axis lines=middle,
  every axis plot/.style={no marks, thick},
  every axis/.style={
    enlarge x limits=false,
    enlarge y limits=auto,
    xlabel style={at={(current axis.right of origin)},anchor=west},
    ylabel style={at={(current axis.above origin)},anchor=south},
    % extra x and y ticks for asymptotes
    extra x tick style={grid=major, grid style={dashed,black}},
    extra y tick style={grid=major, grid style={dashed,black}},
    /pgf/number format/1000 sep={},
  },    
  tick label style = {font=\footnotesize, fill=\currentbgcolor, inner sep=0pt},    
  xticklabel shift=1pt,    
  scaled ticks=false,
  scaled y ticks=false,
  unbounded coords=jump,
  samples=101,
  minor tick length=0,
  scatter/classes={
    open={mark=*,draw=black,fill=white, mark size=2, thin},
    closed={mark=*,draw=black,fill=black, mark size=2, thin}
  },
  width=3in,
}
\pgfplotsset{open/.style={only marks, mark=*, fill=white, thin, mark size=2}}
\pgfplotsset{closed/.style={only marks, mark=*, draw=black, fill=black,
    thin, mark size=2}}
\pgfplotsset{labeled/.style={nodes near coords, only marks, mark=*, draw=black, fill=black, thin, mark size=2, point meta=explicit symbolic}}

\pgfplotsset{gridgraph/.style={clip=false, axis line style={shorten
      >=-8pt}, xlabel style={xshift=8pt}, ylabel style={yshift=8pt},
    enlargelimits=false, grid=major, tick style={black}}} 

\newcommand{\axes}[1][]{
  \begin{tikzpicture}
    \begin{axis}[xtick=\empty, ytick=\empty, ymin=-2, ymax=2,
      xmin=-4.25, xmax=4.25, #1]
    \end{axis}
  \end{tikzpicture}
}
\usepackage{tikz}
\usetikzlibrary{
  matrix,%
  % enables the snake paths
  decorations.pathmorphing,%
  % enables bigger arrows
  decorations.markings,%
  decorations.pathreplacing,%
  decorations.text,%
  calc,%
  patterns,%
  shapes.geometric,%
  arrows,
  decorations.shapes,
  positioning,
  plotmarks,
  intersections
}
\tikzset{>=latex} % sets default arrow type in tikz
\tikzset{font=\normalsize}
\tikzset{mark size=1.5pt, mark options=thin}
\tikzset{pin distance=4pt,
  every pin edge/.style={<-, thin, shorten <= -2pt}}

\interfootnotelinepenalty=10000
\renewcommand{\thefootnote}{(\arabic{footnote})}

\usepackage{multirow, bigdelim}

% change to \usepackage[sols]{handout} to see solutions
\usepackage[sols]{handout}

\newcommand\iscurrentenumi[1]{%
  \ifthenelse{\equal{#1}{\theenumi}}{}{#1}}
\setenumerate[1]{ref=\#\arabic*}
\setenumerate[2]{ref=\protect\iscurrentenumi{\theenumi}(\alph*)}
\newcommand\enumiiprefix[2]{%
  \ifthenelse{{\equal{#1}{\theenumi}}\AND{\equal{#2}{(\alph{enumii})}}}{}{}%
  \ifthenelse{{\equal{#1}{\theenumi}}\AND{\NOT{\equal{#2}{(\alph{enumii})}}}}{#2}{}%
  \ifthenelse{\equal{#1}{\theenumi}}{}{#1#2}%
}
\setenumerate[3]{ref=\protect\enumiiprefix{\theenumi}{(\alph{enumii})}\roman*}

\makeatletter

\newdimen\SOUL@dimen %new
\def\SOUL@ulunderline#1{{%
    \setbox\z@\hbox{#1}%
    % \dimen@=\wd\z@
    \SOUL@dimen=\wd\z@ %new
    \dimen@i=\SOUL@uloverlap
    \advance\SOUL@dimen2\dimen@i %\dimen@ exchanged too
    \rlap{%
      \null
      \kern-\dimen@i
      % \SOUL@ulcolor{\SOUL@ulleaders\hskip\dimen@}%
      \SOUL@ulcolor{\SOUL@ulleaders\hskip\SOUL@dimen}% new
    }%
    \unhcopy\z@
  }}

\makeatother

\newcommand{\related}[1]{%
\fcolorbox{black}{yellow}{%
  %\parbox[t]{12pt}{$\clubsuit$}%
  \parbox[t]{\linewidth-8pt}{\emph{Related problems:} #1}}}

\renewcommand{\answer}[1]{#1}

\definecolor{purple}{rgb}{0.5,0,1.0}
\definecolor{hotpink}{rgb}{1, 0, 0.5}

\definecolor{skyblue}{rgb}{0, 0.5, 1}

\newcommand{\highlight}[1]{\boldmath\hl{\textbf{#1}}\unboldmath}

\newcommand{\goal}[1]{\textbf{\textcolor{skyblue}{#1}}}

\newcommand{\concaveup}{\raisebox{-1pt}{\tikz[x=0.075cm,y=0.075cm] \draw (-2, 4) parabola bend (0, 0) (2, 4);}}
\newcommand{\concavedown}{\raisebox{-1pt}{\tikz[x=0.075cm,y=0.075cm] \draw (-2, -4) parabola bend (0, 0) (2, -4);}}

\usepackage{amsthm}

% make URLs print in the same font
\usepackage{url}
\urlstyle{same}

\usepackage{pifont}
\def\exend{\hfill\ding{118}}

%\makeatletter\@addtoreset{equation}{section}\makeatother
%\renewcommand{\theequation}{\arabic{section}.\arabic{equation}}

\newtheoremstyle{theorem}% name
{8pt}%      Space above
{0pt}%      Space below
{\itshape}%         Body font
{}%         Indent amount (empty = no indent, \parindent = para indent)
{\bfseries}% Thm head font
{.}%        Punctuation after thm head
{.5em}%     Space after thm head: " " = normal interword space;
                                %       \newline = linebreak
{}%         Thm head spec (can be left empty, meaning `normal')

\theoremstyle{theorem}

\let\Theorem\undefined
\newtheorem{Theorem}[equation]{Theorem}
\let\Definition\undefined
\newtheorem{Definition}[equation]{Definition}
\let\Summary\undefined
\newtheorem{Summary}[equation]{Summary}
\let\Conjecture\undefined
\newtheorem{Conjecture}[equation]{Conjecture}
\let\Fact\undefined
\newtheorem{Fact}[equation]{Fact}

\renewcommand{\thesubsection}{\arabic{subsection}}
\makeatletter
\def\@seccntformat#1{\@ifundefined{#1@cntformat}%
   {\csname the#1\endcsname\quad}%       default
   {\csname #1@cntformat\endcsname}}%    enable individual control
\newcommand\section@cntformat{}
\makeatother

  \usepackage{exam}

\usepackage{fancyhdr}
\lhead{\textcolor{white}{This content is protected and may not be shared, uploaded, or distributed.}}
\rhead{}
\rfoot{\vspace*{\fill}\footnotesize\textcopyright{} \emph{Harvard Math Ma}}
\renewcommand{\headrulewidth}{0pt}
\pagestyle{fancy}
\begin{document}

  \title{Exam 2, Practice 1}

\instructions{instructions.tex}{}

\listofproblems

\begin{enumerate}
\item \label{fall2019-midterm-practice1-2-modified} \points{12} 

  Below is the graph of a function $f$ and a table of values for a function $g$.

  \begin{minipage}{0.5\textwidth}
    \begin{tikzpicture}
      \begin{axis}[ xtick={-3, -2, -1, 0, 1, 2, 3, 4, 5}, ytick={-1, 0, 1, 2, 3, 4, 5}, gridgraph]
        \addplot+[open] 
        coordinates {(-1, 3) (1, 5) (1, 3) (2, 4)};%
        \addplot+ [closed] table {
          -3 4
          -1 0
          2 1
          5 5
          1 2
        };%
        \addplot [ domain=1:2, ] {x^2 - 2*x +4};%
        \addplot [ domain=-3:-1, ] {-2*x-2};%
        \addplot [ domain=-1:1, ] {x+4};%
        \addplot [ domain=2:3, ] {1};%
        \addplot [ domain=3:5, ] {2*x-5} node[pos=0.1, right]{$y = f(x)$};%
      \end{axis}
    \end{tikzpicture}
  \end{minipage}
  \begin{minipage}{0.3\textwidth}
    \begin{tabular}{|c|c|r|}
      \hline
      $x$&$g(x)$&$g'(x)$\\[5pt]
      \hline
      0&4&$2$\\[5pt]
      \hline
      1&5&$-3$\\[5pt]
      \hline
      2&3&$7$\\[5pt]
      \hline
      3&2&2\\[5pt]
      \hline
      4&$ 5$&$-1$\\[5pt]
      \hline
      5&1&$-6$\\[5pt]
      \hline
    \end{tabular}
  \end{minipage}

  \begin{enumerate}
  \item
    \begin{problem}[0.25in]
      List all points in the interval $(-3, 5)$ where $f$ is not differentiable.
    \end{problem}

    \begin{solution}
      $f$ is discontinuous at $x = -1, 1, 2$, so it's not differentiable there by {{\#4}} on \href{https://canvas.harvard.edu/courses/138327/files/20921670?wrap=1}{the worksheet ``Continuity''}. Also, the graph has a corner at $x = 3$, so $f$ isn't differentiable there. So, $f$ is not differentiable at $x = \boxed{-1, 1, 2, 3}$. 
    \end{solution}

  \item % added by Janet

    Find the following limits; please show your work / explain your reasoning, and simplify your answers. If a limit fails to exist because it goes to $\infty$ or $-\infty$, indicate which.

    \begin{enumerate}
    \item
      \begin{problem}[\vfill]
        $\displaystyle \lim_{x \to 3} \dfrac{f(x)}{2^{-x}}$
      \end{problem}

      \begin{solution}
        As $x \to 3$, $f(x) \to 1$ and $2^{-x} \to 2^{-3} = \dfrac{1}{8}$. So, $\displaystyle \lim_{x \to 3} \dfrac{f(x)}{2^{-x}} = \frac{1}{1/8} = \boxed{8}$. 
      \end{solution}

    \item
      \begin{problem}[\vfill]
        $\displaystyle \lim_{x \to 1} |f(x) - 4|$
      \end{problem}

      \begin{solution}
        Since $\displaystyle \lim_{x \to 1} f(x)$ does not exist, we need to look at the one-sided limits. As $x \to 1^+$, $f(x) \to 3$, so $|f(x) - 4| \to |3 - 4| = 1$. As $x \to 1^-$, $f(x) \to 5$, so $|f(x) - 4| \to |5 - 4| = 1$. Since $\displaystyle \lim_{x \to 1^+} |f(x) - 4|$ and $\displaystyle \lim_{x \to 1^-} |f(x) - 4|$ are both equal to 1, $\displaystyle \lim_{x \to 1} |f(x) - 4| = \boxed{1}$. 
      \end{solution}

    \item
      \begin{problem}[\vfill\newpage]
        $\displaystyle \lim_{h \to 0} \frac{f(3 + h) - f(3)}{h}$
      \end{problem}

      \begin{solution}
        You should recognize this as the definition of $f'(3)$; since $f$ isn't differentiable at $3$, this limit \fbox{does not exist}.

        \emph{Alternate approach:} Even if you didn't immediately recognize this as the limit definition of the derivative, you can still do this problem! The expression $\dfrac{f(3 + h) - f(3)}{h}$ should jump out as being an average rate of change / slope of a secant line. In particular, it's the slope of the secant line between $(3, f(3))$ and $(3 + h, f(3 + h))$. When $h$ is a little bigger than 0, this slope is 2; when $h$ is a little less than 0, this slope is 0. So, $\displaystyle \lim_{h \to 0^+} \frac{f(3 + h) - f(3)}{h} = 2$ and $\displaystyle \lim_{h \to 0^-} \frac{f(3 + h) - f(3)}{h} = 0$. Since the one-sided limits don't equal each other, $\displaystyle \lim_{h \to 0} \frac{f(3 + h) - f(3)}{h}$ does not exist. 
      \end{solution}

    \item
      \begin{problem}[\vfill]
        $\displaystyle \lim_{x \to -2^+} \frac{3^x}{f(x) - 2}$
      \end{problem}

      \begin{solution}
        As $x \to -2^+$, $3^x \to 3^{-2} = \dfrac{1}{9}$ and $f(x) \to 2$, so $f(x) - 2 \to 0$. Since the denominator approaches 0 and the numerator doesn't, we need to think about the sign of the denominator.

        When $x$ is a little bigger than $-2$ (like $x = -1.99$), $f(x)$ is a little less than 2, so $f(x) - 2$ is a little less than 0. That means $\displaystyle \frac{3^x}{f(x) - 2} = \frac{\text{something close to $\frac{1}{9}$}}{\text{something just a little less than 0}}$, which is a very negative number. So, $\displaystyle \lim_{x \to -2^+} \frac{3^x}{f(x) - 2} = \boxed{-\infty}$.
      \end{solution}

    \end{enumerate}

  \item
    \begin{problem}[\nstretch{2}]
      Let $a(x) = 3 f(x) g(x)$. Find $a'(4)$. 
    \end{problem}

    \begin{solution}
      \begin{align*}
        a'(x) &= 3 [f(x) g(x)]' \text{ by the Constant Multiple Rule} \\
        &= 3 [ f'(x)g(x)+f(x)g'(x) ] \text{ by the Product Rule} \\
        \intertext{Therefore,}
        a'(4) &= 3[ f'(4)g(4)+f(4)g'(4)] \\
        &= 3 [ (2)(5)+(3)(-1)] \\
        &= \boxed{21}
      \end{align*}
    \end{solution}

  \end{enumerate}

  \pspace{\newpage}

\item \points{9} 
 The population of Oz is growing exponentially. On January 1, 2022, the population was 100,000 people. On January 1, 2024, the population was 125,000 people.

\begin{enumerate}
    \item \begin{problem}[\vfill] Write a formula for $P(t)$,  the population of Oz $t$ years after January 1, 2020. \end{problem}
    
    \begin{solution}
        We know that $P(2)=100,000$ and $P(4)=125,000$. Since we know that $P(t)$ is an exponential function, $P(t)=C\cdot b^t$, we also know that $P(2)=C\cdot b^2$ and $P(4)=C\cdot b^4$. We can solve the following system of equations for $C$ and $b$:
        \begin{align*}
            C \cdot b^4 &= 125,000\\
            C\cdot b^2 &= 100,000
        \end{align*}
        Dividing the first equation by the second gives us:
        \begin{align*}
            \dfrac{C\cdot b^4}{C \cdot b^2} &= \frac{125,000}{100,000}\\
            b^2 &= \frac{5}{4}\\
            b &= \left(\frac{5}{4}\right)^{1/2}
        \end{align*}
        Now we can use $b=\left(\frac{5}{4}\right)^{\frac{1}{2}}$ find $C$:
    \begin{align*}
        C \cdot b^2 &= 100,000\\
        C \cdot \left(\left(\frac{5}{4}\right)^{\frac{1}{2}}\right)^2 &= 100,000\\
        C \cdot \frac{5}{4} &= 100,000\\
        C &= 100,000 \cdot \frac{4}{5}\\
        C &= 80,000
    \end{align*}
    So the formula for $P(t)$ is $\boxed{P(t)=80000\cdot \left(\frac{5}{4}\right)^{t/2}}$
    \end{solution}
    \item \begin{problem}Sketch a graph of $P(t)$. Label the y-intercept, any known points, and any asymptotes.\\
    \end{problem}
    \begin{problemonly}
        \axes[]
      \end{problemonly} 

      \begin{solution}
          \begin{tikzpicture}
            \begin{axis}[ yticklabels={$4$}, xtick=\empty, ylabel={population (thousands)}, xlabel={$t$ (years after Jan 1, 2020)}, ]
              \addplot+[domain=-1:6] {80*1.25^(.5*x)};
              \filldraw[gray, below] (2,100) circle (2pt) node{ $(2,100000)$};
              \filldraw[gray, below] (4,125) circle (2pt) node{ $(4, 125000)$};
              
              \node[above left, outer sep=2pt] (0,80) {80000};
              
            \end{axis}
          \end{tikzpicture}
      \end{solution}
    \item \begin{problem}[\vfill] Do you expect that the the population of Oz be greater than 150,000, less than 150,000, or equal to 150,000 in 2026? Explain your reasoning in 1-2 sentences.\\
    \end{problem}
    \begin{solution}
        We expect that the population will be greater than 150,000. Exponential functions grow at a rate proportional to their value. As the population increases in size, so will its growth rate. We can also see this from the concave up, increasing nature of the function's graph.
    \end{solution}
    \item \begin{problem}$P'(t)$ is a(n) \textbf{increasing / decreasing} (Circle one) function.\\
    \end{problem}
    \begin{solution} 
    $P'(t)$ is an {\bf increasing }function. Concave up functions have increasing derivatives. We can also sketch a couple of tangent lines onto the function and see that the slopes of the tangent lines are increasing as $t$ increases.
    \end{solution}
    \item \begin{problem}$P'(t)$ is \textbf{concave up / concave down} (Circle one).\\
    \end{problem}
    \begin{solution}
        Exponential functions' derivatives are proportionate to themselves. This means that the derivative will be a multiple of the original function. Looking at $P(t)$, we know the derivative is positive ($P$ increases everywhere), is increasing, and is also an exponential. This means the derivative will also take the form of exponential growth, so it will also be {\bf concave up,} just like $P(t)$.
    \end{solution}
    \newpage
\end{enumerate}

\item \points{8} 

  \begin{enumerate}
  \item In each part, either sketch an example of a function $f$ with the given properties, or say that no such function exists. Make sure any sketches clearly meet the given the conditions.

    \begin{enumerate} 

    \item \label{fall2017-prefinal-3b}

      \begin{problem}
        $f$ has domain $[-2, 2)$, is continuous on its domain, and has a global maximum value\probonly{\footnote{This is referring to the $y$-value.}} of 3 and no global minimum on $[-2, 2)$.
      \end{problem}

      \begin{problemonly}
        \axes[xmin=-2, xmax=2, ymin=-1, ymax=3, gridgraph, xtick={-2, -1, 1, 2}, ytick={-1, 1, 2, 3}, axis equal image]
      \end{problemonly}

      \begin{solution}
        Here are two examples:
        \begin{center}
          \answer{
            \begin{tikzpicture}
              \begin{axis}[xmin=-2, xmax=2, ymin=-1, ymax=3, gridgraph, axis equal image]
                \addplot+[sharp plot] coordinates {(-2, 3) (2, 0)};%
                \addplot+[closed] coordinates { (-2, 3)};%
                \addplot+[open] coordinates {(2, 0)};
              \end{axis}
            \end{tikzpicture}
            \hspace{0.5in}
            \begin{tikzpicture}
              \begin{axis}[xmin=-2, xmax=2, ymin=-1, ymax=3, gridgraph, axis equal image]
                \addplot+[sharp plot] coordinates{ (-2, 1) (-1, 3) (0, 3) (1, 1) (2, -1)}; %
                \addplot+[closed] coordinates {(-2, 1)}; %
                \addplot+[open] coordinates {(2, -1)}; %
              \end{axis}
            \end{tikzpicture}
          }
        \end{center}
      \end{solution}

    \item \label{fall2017-midterm-1b}
      \begin{problem}
        $f$ is continuous on $(-\infty, \infty)$, and $\displaystyle \lim_{x \to 3} f(x)$ does not exist.
      \end{problem}

      \begin{problemonly}
        \axes[]
      \end{problemonly}       

      \begin{solution}
        \fbox{No such function exists.} Saying $f$ is continuous at 3 means $\displaystyle \lim_{x \to 3} f(x) = f(3)$, so $\displaystyle \lim_{x \to 3} f(x)$ exists. 
      \end{solution}

    \end{enumerate} 

  \item
    \begin{problem}[\vfill\newpage]
      \textbf{Write a formula} for a function $f$ that has $\displaystyle \lim_{x \to 3} f(x) = -\infty$, $\displaystyle \lim_{x \to -\infty} f(x) = 4$, and $\displaystyle \lim_{x \to \infty} f(x) = 4$. (Sketching an example may help!)
    \end{problem}

    \begin{solution}
      Here's a sketch that matches the description:
      \begin{center}
        \begin{tikzpicture}
          \begin{axis}[ymin=-6, ymax=6, xtick={3}, ytick={4}]
            \addplot+[domain=-1:7] {4-3/(x-3)^2}; %
            \addplot+[dashed, domain=-1:7] {4}; %
            \addplot+[dashed, domain=-6:6] (3, x); % 
          \end{axis}
        \end{tikzpicture}
      \end{center}
      \begin{noanswer}
        We can get a function of this shape by starting with $y = \dfrac{1}{x^2}$, shifting it 3 units to the right (which has formula $y = \dfrac{1}{(x - 3)^2}$), reflecting it over the $x$-axis (which has formula $y = - \dfrac{1}{(x - 3)^2}$), and moving it up 4. That gives $\boxed{f(x)=-\frac{1}{(x-3)^2}+4}$.
      \end{noanswer}
      \answeronly{One possibility is $f(x)=-\dfrac{1}{(x-3)^2}+4$.}
    \end{solution}

  \end{enumerate}

\item \label{fall2018-prefinal-practice1-5} \points{10}

  In this problem, you'll use a tangent line approximation to approximate $\sqrt[3]{8.5}$. 

  \solonly{\related{\href{https://canvas.harvard.edu/courses/138327/files/21129250?wrap=1}{Problem Set 18}, {{{{\#4}}(a)}}}}

  \begin{enumerate}
  \item
    \begin{problem}[1in]
      What function will you use to approximate $\sqrt[3]{8.5}$?
    \end{problem}

    \begin{solution} 
      \answer{$f(x) = \sqrt[3]{x}$}
    \end{solution} 

  \item
    \begin{problem}[\vfill\newpage]
      Sketch the graph of the function you chose. (Make your graph large because you will add to it later.) 
    \end{problem}

    \begin{solution}
      You might not remember this graph off the top of your head, but  but you should be able to plot a few points: $(0, 0)$, $(1, 1)$, $(8, 2)$, and $(27, 3)$ are examples of points on the graph. So, that gives us the following graph:
      \begin{center}
        \answer{
          \begin{tikzpicture}
            \begin{axis}[xtick=\empty, ytick=\empty]
              \addplot[domain=0:1] {x^(1/3)};
              \addplot+[domain=1:10] {x^(1/3)}; 
            \end{axis}
          \end{tikzpicture}
        }
      \end{center}      
    \end{solution}

  \item
    \begin{problem}[\vfill]
      Approximate $\sqrt[3]{8.5}$ using a tangent line approximation.
    \end{problem}

    \begin{solution}
      We need to pick a value $x = a$ where we'll find a tangent line. We're trying to approximate $f(8.5)$, so we want to pick a value of $a$ that's close to $8.5$ and where $f(a)$ is easy to evaluate. The closest suitable value is $8$, since we can easily evaluate $f(8) = 2$. 

      The slope of the tangent line at $x = 8$ is $f'(8)$. By the Power Rule, $f'(x) = \dfrac{1}{3} x^{-2/3}$, so $f'(8) = \dfrac{1}{3} 8^{-2/3} = \dfrac{1}{3} \cdot \dfrac{1}{8^{2/3}} = \dfrac{1}{3} \cdot \dfrac{1}{(8^{1/3})^2} = \dfrac{1}{3} \cdot \dfrac{1}{2^2} = \dfrac{1}{12}$. We know $(8, 2)$ is a point on the tangent line, so the tangent line has equation $y - 2 = \dfrac{1}{12} (x - 8)$, or $y = \dfrac{1}{12} (x - 8) + 2$. This should be a good approximation for $f(x)$ when $x$ is close to 8. That is,
      \begin{align*}
        \sqrt[3]{x} & \approx \dfrac{1}{12} (x - 8) + 2 \text{ for $x$ near 8} \\
        \intertext{Plugging in $x = 8.5$,}
        \sqrt[3]{8.5} & \approx \dfrac{1}{12} \cdot \dfrac{1}{2} + 2 \\
        &= \boxed{\frac{49}{24}}
      \end{align*}
    \end{solution} 

  \item
    \begin{problem}[\nstretch{0.75}]
      Is your approximation an upper bound or lower bound for the actual value of $\sqrt[3]{8.5}$? Explain, using your picture in (b) to illustrate. 
    \end{problem}

    \begin{solution}
      Here's a picture to illustrate; we found the tangent line to $y = \sqrt[3]{x}$ at $x = 8$:
      \begin{center}
        \begin{tikzpicture}
          \begin{axis}[xtick={8}, ytick=\empty]
            \addplot[domain=0:1] {x^(1/3)};%
            \addplot+[domain=1:10] {x^(1/3)};%
            \addplot+[hotpink, domain=4:10] {(x - 8) / 12 + 2}; %            
          \end{axis}
        \end{tikzpicture}
      \end{center}      
      Our approximation is the height of the tangent line at $x = 8.5$, which is higher than the height of the graph of $\sqrt[3]{x}$ at $x = 8.5$. Therefore, our approximation is an \fbox{upper bound}. 
    \end{solution}

  \end{enumerate}

\item \points{24}  

  Let $f(x) = \dfrac{x^2 - 4}{3x^2 + 1}$.

  \begin{enumerate}

  \item
    \begin{problem}[\vfill]
      Is $f$ even, odd, or neither? Justify your answer. 
    \end{problem}

    \begin{solution}
      As always, we look at $f(-x)$:
      \begin{align*}
        f(-x) &= \dfrac{(-x)^2 - 4}{3(-x)^2 + 1} \\
        &= \dfrac{x^2 - 4}{3x^2 + 1}
      \end{align*}
      This is equal to $f(x)$, so $f$ is \fbox{even}. 
    \end{solution}

  \item
    \begin{problem}[\vfill]
      Find and classify all discontinuities of $f$ (that is, say whether each discontinuity is a removable discontinuity, jump discontinuity, or vertical asymptote). Justify your classifications using limits.
    \end{problem}

    \begin{solution}
      Since $f$ is given by a single (non-piecewise) formula, its discontinuities are exactly where it is undefined. $f(x)$ is undefined when its denominator is 0:
      \begin{align*}
        3x^2 + 1&= 0 \\
        x^2 &= - \frac{1}{3}
      \end{align*}
      But this is impossible, so $f(x)$ is never undefined. Therefore, $f$ has \fbox{no discontinuities}.
    \end{solution}

  \item
    \begin{problem}[\vfill\newpage]
      Find all horizontal asymptotes of $f$. Justify your answer completely. 
    \end{problem}

    \begin{solution}
      To find the horizontal asymptotes, we must evaluate $\displaystyle \lim_{x \to \infty} f(x)$ and $\displaystyle \lim_{x \to -\infty} f(x)$.

      Let's evaluate $\displaystyle \lim_{x \to \infty} f(x)$ first. As $x \to \infty$, both the numerator and denominator have a term that $\to \infty$. As we've seen, in a sum where some terms go to $\pm \infty$, we only need to keep the fastest-growing term.\footnote{Why does this work? In the sum $x^2 - 4$, when $x$ gets very big, the $-4$ is tiny compared to $x^2$, so $x^2 - 4 \approx x^2$ when $x$ is very big. Similarly, when $x$ is very big, $3x^2 + 1 \approx 3x^2$.} So, \[ \lim_{x \to \infty} \frac{x^2 - 4}{3x^2 + 1} = \lim_{x \to \infty} \frac{x^2}{3x^2} = \lim_{x \to \infty} \frac{1}{3} = \frac{1}{3}.\]
      Similarly,
      \[ \lim_{x \to -\infty} \frac{x^2 - 4}{3x^2 + 1} = \lim_{x \to -\infty} \frac{x^2}{3x^2} = \lim_{x \to -\infty} \frac{1}{3} = \frac{1}{3}.\]
      So, the only horizontal asymptote is $\boxed{y = \frac{1}{3}}$. 
    \end{solution}

  \item
    \begin{problem}[\vfill]
      On what intervals is $f(x)$ increasing? On what intervals is it decreasing?
    \end{problem}

    \begin{solution}
      To answer this, let's make a sign chart for $f'$. Using the Quotient Rule,
      \begin{align*}
        f'(x) &= \dfrac{(3x^2 + 1) 2x - (x^2 - 4)6x}{(3x^2 + 1)^2} \\
        &= \dfrac{6x^3 + 2x - 6x^3 + 24x}{(3x^2 + 1)^2} \\
        &= \dfrac{26 x}{(3x^2 + 1)^2}
      \end{align*}
      To make our sign chart, we start by finding where $f'(x)$ is 0 or undefined.
      \begin{alignat*}{4}
        f'(x) &= 0 & \qquad \text{ or } \qquad && f'(x) &= \text{undefined} \\
        26x &= 0 & && (3x^2 + 1)^2 &= 0 \\
        x &= 0 &&& 3x^2 + 1 &= 0 \\
        &&&& x^2 &= - \frac{1}{3} \\        
        & &&& \multispan3{\hspace{0.4in}Impossible!}
      \end{alignat*}
      So, the only value we need to mark on our sign chart is $x = 0$.

      Then, we test points on either side of $0$. For example, $f'(1) = \dfrac{26}{4^2} > 0$ and $f'(-1) = \dfrac{-26}{4^2} < 0$. So, here's the sign chart for $f'$:
      \begin{center}
        \begin{tikzpicture}
          \begin{axis}[axis y line=none, x axis line style={-}, xtick={0}, ymin=-2, ymax=2, height=1.25in, width=3in, clip=false]
            \addplot+[domain=-0.2:0.1, very thin]{0};
            \draw(-0.2, 0) node[anchor=south west] {$f$};%
            \draw (-0.2, -1.5) node[right] {sign of $f'$};%
            \draw (-0.05, -1.5) node{$-$};%
            \draw[->, xshift=2pt] (-0.0933333, 1.5) -- (-0.00666667, 0.5);%
            \draw (0.05, -1.5) node{$+$};%
            \draw[->, xshift=2pt] (0.00666667, 0.5) -- (0.0933333, 1.5);%
            \draw[->, xshift=2pt] (-0.00666667, 0.5) -- (0.00666667, 0.5);%
          \end{axis}
        \end{tikzpicture}
      \end{center}
      So, \fbox{$f$ is decreasing on $(-\infty, 0)$ and increasing on $(0, \infty)$}.
    \end{solution}

  \end{enumerate}
  It turns out that $f''(x) = \dfrac{26 (1 - 9x^2)}{(3x^2 + 1)^3}$. 
  \begin{enumerate}[resume]
  \item
    \begin{problem}[\vfill]
      On what intervals is $f(x)$ concave up? On what intervals is it concave down?
    \end{problem}

    \begin{solution}
      To answer this, let's make a sign chart for $f''$. As usual, we start by finding where $f''(x)$ is 0 or undefined. We're told that $f''(x) = \dfrac{26 (1 - 9x^2)}{(3x^2 + 1)^3}$.            
      \begin{alignat*}{4}
        f'(x) &= 0 & \qquad \text{ or } \qquad && f'(x) &= \text{undefined} \\
        1 - 9x^2 &= 0 &&& (3x^2 + 1)^3 & = 0 \\
        x^2 &= \frac{1}{9} &&& 3x^2 + 1 &= 0 \\
        x &= \pm \frac{1}{3} &&& x^2 &= - \frac{1}{3} \\
        & &&& \multispan3{\hspace{0.4in}Impossible!}
      \end{alignat*}
      So, the values we need to mark on our sign chart are $\pm \dfrac{1}{3}$. Testing points gives us:
      \begin{center}
        \begin{tikzpicture}
          \begin{axis}[axis y line=none, x axis line style={-}, xtick={-0.333333, 0.333333}, xticklabels={$-\frac{1}{3}$, $\frac{1}{3}$}, ymin=-2, height=1in, width=3in, clip=false]
            \addplot+[domain=-1.66667:1, very thin]{0};
            \draw (-1.66667, -1.5) node[right] {sign of $f''$};
            \draw (-0.666667, -1.5) node{$-$};%
            \draw (0, -1.5) node{$+$};%
            \draw (0.666667, -1.5) node{$-$};%
            \draw (-1.66667, 0) node[anchor=south west] {$f$};
            \draw (-0.666667, 1) node{\concavedown};%
            \draw (0, 1) node{\concaveup};%
            \draw (0.666667, 1) node{\concavedown};%
          \end{axis}
        \end{tikzpicture}
      \end{center}
      For example, $f''(-1) = \dfrac{26 (-8)}{4^3} < 0$, $f''(0) = \dfrac{26}{1} > 0$, and $f''(1) = \dfrac{26 (-8)}{4^3} < 0$.

      So, \fbox{$f$ is concave down on $\left(-\infty, -\frac{1}{3} \right) \cup \left( \frac{1}{3}, \infty \right)$ and concave up on $\left( - \frac{1}{3}, \frac{1}{3} \right)$}.
    \end{solution}

  \item
    \begin{problem}[0.25in]
      What are the turning points of $f$?
    \end{problem}

    \begin{solution}
      The only place where $f$ switches from increasing to decreasing or vice versa is $x = \boxed{0}$. 
    \end{solution}

  \item

    \begin{problem}[0.25in]
      What are the inflection points of $f$?  
    \end{problem}

    \begin{solution}
      The concavity of $f(x)$ changes at $x = \boxed{ - \frac{1}{3}, \frac{1}{3}}$. 
    \end{solution}

    \pspace{\newpage}

  \item
    \begin{problem}[1in]
      What are the $x$- and $y$-intercepts of $f$?
    \end{problem}

    \begin{solution}
      The \answer{$y$-intercept is $f(0) = -4$}. To find the $x$-intercepts, we solve $f(x) = 0$:
      \begin{align*}
        \dfrac{x^2 - 4}{3x^2 + 1} &= 0 \\
        x^2 - 4 &= 0 \\
        (x - 2)(x + 2) &= 0 \\
        x &= 2, -2
      \end{align*}
      So, the \answer{$x$-intercepts are $x = \pm 2$}. 
    \end{solution}

  \item
    \begin{problem}[\newpage]
      Sketch a rough graph of $f(x)$ that incorporates all of the above
      information.
    \end{problem}

    \begin{solution}
      As a first draft, let's mark the intercepts we found, as well as the fact that both $\displaystyle \lim_{x \to \infty} f(x)$ and $\displaystyle \lim_{x \to -\infty} f(x)$ are $\dfrac{1}{3}$:
      \begin{center}
        \begin{tikzpicture}
          \begin{axis}[xtick={-2, 2}, ytick={-4, 0.333}, yticklabels={$-4$, $\frac{1}{3}$}, width=4in, height=2.5in]
            \addplot+[closed] coordinates {(0, -4) (-2, 0) (2, 0) };
            \addplot+[dashed, domain=3:4] {1/3};
            \addplot+[dashed, domain=-4:-3] {1/3}; 
          \end{axis}
        \end{tikzpicture}
      \end{center}
      Next, let's do a draft in light blue that shows where $f$ is increasing and decreasing but ignores concavity:
      \begin{center}
        \begin{tikzpicture}
          \begin{axis}[xtick={-2, 2}, ytick={-4, 0.333}, yticklabels={$-4$, $\frac{1}{3}$}, width=4in, height=2.5in]
            \addplot+[closed] coordinates {(0, -4) (-2, 0) (2, 0) };
            \addplot+[dashed, domain=3:4] {1/3};
            \addplot+[dashed, domain=-4:-3] {1/3};
            \addplot+[skyblue, dotted, domain=-4:-2] { (x^2 - 4) / (3*x^2 + 1)}; %
            \addplot+[skyblue, dotted, domain=2:4] { (x^2 - 4) / (3*x^2 + 1)}; %
            \addplot[skyblue, dotted, sharp plot] coordinates { (-2, 0) (0, -4) (2, 0) }; %
          \end{axis}
        \end{tikzpicture}
      \end{center}
      Finally, let's add in the concavity and the fact that the function is even (so the graph is symmetric over the $y$-axis). To make the concavity clear, I'll draw the concave up portions in green and the concave down portions in pink:
      \begin{center}
        \begin{tikzpicture}
          \begin{axis}[xtick={-2, -0.333, 0.333, 2}, xticklabels={$-2$, $-\frac{1}{3}$, $\frac{1}{3}$, $2$}, ytick={-4, 0.333}, yticklabels={$-4$, $\frac{1}{3}$}, width=4in, height=2.5in]
            \addplot+[domain=-4:-1/3, green!70!black] { (x^2 - 4) / (3*x^2 + 1)};
            \addplot+[domain=-1/3:1/3, hotpink] { (x^2 - 4) / (3*x^2 + 1)};
            \addplot+[domain=1/3:4, green!70!black] { (x^2 - 4) / (3*x^2 + 1)};
            \addplot+[dashed, domain=3:4] {1/3};
            \addplot+[dashed, domain=-4:-3] {1/3};
          \end{axis}
        \end{tikzpicture}
      \end{center}
      So, here's our final graph:
      \begin{center}
        \answer{
          \begin{tikzpicture}
            \begin{axis}[xtick={-2, -0.333, 0.333, 2}, xticklabels={$-2$, $-\frac{1}{3}$, $\frac{1}{3}$, $2$}, ytick={-4, 0.333}, yticklabels={$-4$, $\frac{1}{3}$}, width=4in, height=2.5in]
              \addplot+[domain=-4:4] { (x^2 - 4) / (3*x^2 + 1)};
              \addplot+[dashed, domain=-4:4] {1/3}; 
            \end{axis}
          \end{tikzpicture}
        }
      \end{center}
    \end{solution}

    \begin{grading}
      The main thing we're looking for is that your graph is consistent with all of the earlier parts of the problem. So, when you're done with the problem, go back and check that your graph really fits with everything you said earlier!
    \end{grading}

  \end{enumerate}

\item \label{fall2013-1a-final-review2-04} \points{12} 

  \begin{problem}[\newpage]
    Christy wants to build a big rectangular pool in her backyard, but she also wants to keep as much of the pretty green grass as possible. Because of this, she has the following three requirements:
    \begin{itemize}[nosep]
    \item The pool must have an area of $500$ square meters for the water.
    \item The pool must have a rectangular concrete border that is at least $1$ meter wide.
    \item She wants to minimize the amount of area that the pool and concrete border take up.
    \end{itemize}
    What should the dimensions of her pool be to achieve these three requirements?
  \end{problem}

  \begin{solution}
    \highlight{Goal:} Minimize \goal{total area} (water plus border)

    Just using common sense, if Christy wants the pool to take up the least amount of area, she should make the concrete borders exactly 1 meter wide (because that's the thinnest possible). Let's draw two possible pools:

    \begin{tikzpicture}[x=0.75cm, y=0.75cm]
      \draw (0, 0) -- (8, 0) -- (8, 6) -- (0, 6) -- (0, 0);
      \draw[fill=black!10!white] (1, 1) -- (7, 1) -- (7, 5) -- (1, 5) -- (1, 1);
      \draw (4, 3) node{water};
      \draw (4, 5.5) node{border};
      \draw[decorate, decoration={brace, mirror}, red] (1, 1) --
      node[right, xshift=1pt, red]{$y$} (1, 5);
      \draw[decorate, decoration={brace}, red] (1, 1) --
      node[above, yshift=2pt, red]{$x$} (7, 1);
      \draw[|-|] (5, 0) -- node[fill=white]{1} (5, 1);
      \draw[|-|] (5, 5) -- node[fill=white]{1} (5, 6);
      \draw[|-|] (0, 3) -- node[fill=white]{1} (1, 3);
      \draw[|-|] (7, 3) -- node[fill=white]{1} (8, 3);
    \end{tikzpicture}
    \qquad 
    \begin{tikzpicture}[x=0.75cm, y=0.75cm]
      \draw (0, 0) -- (11.6, 0) -- (11.6, 4.5) -- (0, 4.5) -- (0, 0);
      \draw[fill=black!10!white] (1, 1) -- (10.6, 1) -- (10.6, 3.5) -- (1, 3.5) -- (1, 1);
      \draw (5.8, 2.25) node{water};
      \draw (5.8, 4) node{border};
      \draw[decorate, decoration={brace, mirror}, red] (1, 1) --
      node[right, xshift=1pt, red]{$y$} (1, 3.5);
      \draw[decorate, decoration={brace}, red] (1, 1) --
      node[above, yshift=2pt, red]{$x$} (10.6, 1);
      \draw[|-|] (6.8, 0) -- node[fill=white]{1} (6.8, 1);
      \draw[|-|] (6.8, 3.5) -- node[fill=white]{1} (6.8, 4.5);
      \draw[|-|] (0, 2.25) -- node[fill=white]{1} (1, 2.25);
      \draw[|-|] (10.6, 2.25) -- node[fill=white]{1} (11.6, 2.25);
    \end{tikzpicture}

    \highlight{Define variables:}
    \begin{itemize}[nosep]
    \item \goal{$A$ = total area} (water plus border)
    \item $x, y$ = lengths labeled above
    \end{itemize}

    \highlight{Formula for goal:} The width of the water plus border is $x + 2$, and the length is $y + 2$, so
    \[ \goal{$A$} = (x + 2)(y + 2)\] This formula involves both $x$ and $y$, so we need to \highlight{relate $x$ and $y$ using something else we're told in the problem}.
    \begin{align}
      \text{area of water} &= 500 \text{ m}^2 \nonumber \\
      xy &= 500 \nonumber \\
      y &= \frac{500}{x} \label{fall2013-1a-final-review2-04-relation}
    \end{align}
    Plugging this into our formula for $A$ gives $A$ as a function of just
    $x$,
    \[ \goal{$A(x)$} = (x + 2)\left( \frac{500}{x} + 2 \right) = 500 + \frac{1000}{x} + 2x + 4 = 2x + \frac{1000}{x} + 504.\]
    \highlight{Critical points:} $\displaystyle A'(x) = 2 - 1000 x^{-2} = 2 - \frac{1000}{x^2}$ is undefined when $x = 0$, and it's 0 when
    \begin{align*}
      2 &= \frac{1000}{x^2} \\
      x^2 &= 500 \\
      x &= \pm \sqrt{500} = \pm 10 \sqrt{5}
    \end{align*}
    But the only sensible values of $x$ are $x > 0$, so the only critical point in the domain is $x = 10 \sqrt{5}$.

    \highlight{Global minimum:} Since there's only one critical point in the domain, we can try the First or Second Derivative Test. Let's try the Second Derivative Test, since it's easy to calculate that $A''(x) = 2000 x^{-3} = \frac{2000}{x^3}$. Then, $A''(10 \sqrt{5}) = \frac{2000}{(10 \sqrt{5})^3} > 0$, so $A(x)$ has a local minimum at $x = 10 \sqrt{5}$. Since this was the only critical point in the domain, this must be the global minimum.

    \highlight{Answer the question:} We can plug our value of $x$ back into (\ref{fall2013-1a-final-review2-04-relation}) to find that the corresponding value of $y$ is $y = 10 \sqrt{5}$. So, \fbox{the water should be a square with sides of length $10 \sqrt{5}$}.
  \end{solution}

\item \label{fall2014-1a-midterm1-1} \points{5} \finishexam 

  Craig just bought a piece of property in Cambridge, and he'd like to build a house on it. He hires a team of architects to design the house for him. They tell him that the cost of building a house whose area is $A$ square feet is $C(A)$ dollars.
  \begin{enumerate}
  \item
    \begin{problem}[\vfill]
      What are the units of $\frac{dC}{dA}$?
    \end{problem}

    \begin{solution}
      The units of $\frac{dC}{dA}$ are the units of $C$ (dollars) over
      the units of $A$ (square feet), or \fbox{dollars per square
        foot}.

      If you prefer to think graphically, remember that we visualize the derivative as the slope of a line tangent to the graph of $C$. Even though we don't know what the graph of $C$ looks like, we know that the axes should be:
      \begin{center}
        \begin{tikzpicture}
          \begin{axis}[xmin=0, xmax=30, ymin=0, ymax=20, ytick=\empty,
            xtick=\empty, xlabel=$A$ (square feet), ylabel=$C$ (dollars),
            width=1.5in, enlarge y limits=false]
            \addplot+[very thin] {0};
          \end{axis}
        \end{tikzpicture}
      \end{center}
      So, a slope on this graph should have units of dollars per square
      foot. 
    \end{solution}

  \item
    If $C'(1000) = 100$, which is the most reasonable conclusion?  Please
    explain your choice. 
    \begin{altenumerate}[label=\protect\maybecircle{\Alph*.}]
    \plainitem It will take about $100$ days to build a $1000$ square foot
      home.
    \plainitem If Craig wants a 1000 square foot home, it will cost about
      \$100 per square foot.
    \plainitem If Craig opts for a 1020 square foot home rather than a
      1000 square foot home, he can expect to pay about \$100 more.
    \circleitem If Craig opts for a 980 square foot home instead of a 1000
      square foot home, he can expect to pay about \$2,000 less. 
    \plainitem If Craig opts for a 1500 square foot home rather than a
      1000 square foot home, he can expect to pay about \$50,000 more.
    \end{altenumerate}

    \pspace{\vfill\newpage}

    \begin{solution}
      The fact that $C'(1000) = 100$ means that, when $A = 1000$ square feet, the instantaneous rate of change in the cost of a house is $100$ dollars per square foot. This means that, if Craig changes the size of the house by $s$ square feet, he should expect the cost to change by $(100 \text{ dollars / square foot})(s \text{ square feet}) = 100s \text{ dollars}$, assuming $s$ is relatively small (because this is only an instantaneous rate). Therefore, choice \answer{D} is the best: a change of $-20$ square feet should make the cost change by about $-2000$ dollars.
    \end{solution}

  \end{enumerate}

\end{enumerate}

\end{document}
